 
The development of temperaments and the ongoing tug-of-war between just tuning and transposability over many centuries is a fascinating aspect in the history of music theory and practice. At the beginning of the 20$^{\rm th}$ century it seemed as if the universal acceptance of the ET had finally settled this issue and, in fact, still today this is the prevailing view. In contrast, we share the opinion that the quest for a better intonation is not over yet and that the ET is probably only an intermediate rather than a final solution.

Looking back at the past 150 years it seems that the search for a better intonation oscillates between enthusiasm and disillusionment. For example, starting with the seminal work by Helmholtz~\cite{helmholtz} in 1865, who was among the first to provide a scientific basis for the sensation of music, many theorists and instrument makers at the end of the 19$^{\rm th}$ century were inspired by the challenge to construct a ``\textit{Reininstrument}'', but the solutions were simply too complicated to become accepted on a broad scale. Then in the early 20$^{\rm th}$ century the interest in just intonation abated, probably in part because many composers were writing music that was highly dissonant and even atonal.

In the second half of the last century, a  renewed interest in just intonation and different ways of tuning arose alongside with an increasing attention on historical performance practices~\cite{DuffinArticle}. The new technologies becoming available constituted  another promoting factor for this process. Following the visionary contributions by Eivind Grove~\cite{GrovenMax}, the emerging computer technology led to a variety of proposals, patents, and software packages which reflect the technological capabilities of the respective time. Unfortunately, apart from few exceptions, none of these approaches reached a broader dissemination, partly because the whole issue had maintained the reputation of being exotic and academic, linked to the microtonal community where 12 tones per octave are considered merely as an exception rather than a rule. 

Meanwhile, however, an ordinary mobile phone has become more powerful than a supercomputer in the 80's, offering new and previously undreamt-of possibilities. For example, solving a system of equations in real time, as proposed in the present work, would have been inconceivable two decades ago. Moreover, digital information technology continues to change the musical landscape and the art of instrument making entirely. The purpose of this project is to demonstrate that by now we have completely different means at our disposal which allow us to consider different approaches and to make dynamic tuning schemes suitable for everyday use.  A systematic evaluation of the tuning method is beyond the scope of this article, but we invite the interested readers to form their own subjective conclusions by testing the freely available software or by listening to the available sound examples on our website \href{http://www.just-intonation.org}{\texttt{just-intonation.org}}.

Finally, electronic communication increasingly enhances the interaction between different musical cultures. On the one hand, it is obvious that many traditional intonation systems throughout the world are increasingly influenced (if not even destroyed) by the Western ET. On the other hand, it should not be underestimated that this interaction also influences the Western world, and it cannot be ruled out that at some point in the future it may become fashionable to deviate from the ET. In addition, it is to be expected that the art of instrument making will continue to evolve rapidly and that on the long run the importance of statically tuned temperaments may decrease. All this suggests that dynamically adapting tuning scheme might become more important in the future. This does not necessarily mean that just intervals are the ultimate goal, and in fact it has been shown by~\cite{mathews} that small deviations from rational frequency ratios may certainly perceived as pleasant, but perhaps there will be a growing interest in more consonant intonation schemes. With our contribution we would like to point out that there is a lot of open space for further research and development in this direction.

\vspace{2mm}
\noindent
\textbf{Acknowledgements}\\
We would like to thank the anonymous reviewers for their substantial reports and in particular for making us aware of the work by John de Laubenfels. We also thank R. W. Duffin for the stimulating exchange of ideas, in particular for his suggestion to consider variable interval sizes. We are also grateful to U. Konrad, who made it possible to record the samples for the application. Finally, we would like to thank A. Whisnant for critical reading of the manuscript. 
